\documentclass[11pt]{article}
\usepackage{fullpage}
\usepackage{times}
\usepackage{hyperref}
\usepackage{amsmath, amssymb}
\usepackage{parskip}

% ─── Title Block ──────────────────────────────────────────────────────────────
\title{%
  \textbf{Interpretable N-Body Physics with GKANs}\\[0.4em]
  \large CPSC 440 Machine Learning --- Project Proposal
}
\author{
  Hanson Sun \quad Eitan Klinger \quad Author Three\\
  \small\texttt{\{jsun33, ekling01, student3\}@cs.ubc.ca}
}

\begin{document}
\maketitle
\thispagestyle{empty}

% ─── 1. Problem Statement ─────────────────────────────────────────────────────
\section*{1. Problem Statement}

% What specific problem are you tackling?
% Introduce any non-ML background that a CS/stat/ECE grad student might not know.
% You do NOT need to introduce ML background covered in the course.

[Describe the problem you are focusing on. Be precise: What is the input? What
is the output or goal? Why is this problem interesting or important?
Introduce any domain-specific background required to understand the setting.]

Approaches towards extracting mathematical equations for n-body systems, where multiple masses affect each other gravitationally, rely on black-box models. Doing so hinders interpretability and limits generalizations.

... more background on MLP approaches?

... stuff about how KAN allows symbolic learning directly without requiring another step

... overall the input would be n-body simulation data, output would be interpretable, symbolic representations of the math governing the simulations...


% ─── 2. Proposed Approach ────────────────────────────────────────────────────
\section*{2. Proposed Approach}

% What do you plan to do?
% Which methods will you apply, extend, or implement?
% What data will you use? How will you obtain / preprocess it?

[Describe your plan concretely. For example:
\begin{itemize}
  \item \textbf{Methods:} Which ML techniques will you use or build upon?
        (e.g., logistic regression, kernel methods, deep nets, graphical models)
  \item \textbf{Dataset:} What data will you work with, and how will you
        acquire / preprocess it?
  \item \textbf{Experimental setup:} How will you evaluate your approach?
        What baselines will you compare against?
\end{itemize}]


% ─── 3. Expected Contribution ────────────────────────────────────────────────
\section*{3. Expected Contribution}

% What will doing this project add to the world?
% Choose one of the project types below and explain your specific contribution.
%
% Common project types (pick one or mix):
%   Application bake-off | New application | Scaling up | Improving performance
%   Generalization to new setting | Survey/perspective | Coding project
%   Reproducibility report | Theory

[State clearly what the contribution of your project will be. Examples:
\begin{itemize}
  \item Showing that method X outperforms Y on application Z, with an
        explanation of why.
  \item A new algorithm that scales method X to datasets of size $n \gg 10^6$.
  \item A structured survey of topic X with original insights on future
        directions.
  \item Negative result: a careful argument for why approach X does not work
        in setting Y, and why.
\end{itemize}]


% ─── 4. Team \& Responsibilities (delete if solo) ────────────────────────────
\section*{4. Team \& Division of Responsibilities}

\begin{center}
\begin{tabular}{lll}
\hline
\textbf{Name} & \textbf{Student \#} & \textbf{Primary Responsibilities} \\
\hline
Author One   & 12345678 & [e.g., data preprocessing, experiments] \\
Author Two   & 87654321 & [e.g., method implementation, writeup]   \\
Author Three & 11223344 & [e.g., baseline comparisons, evaluation] \\
\hline
\end{tabular}
\end{center}


% ─── 5. Timeline ─────────────────────────────────────────────────────────────
\section*{5. Timeline}

\begin{center}
\begin{tabular}{ll}
\hline
\textbf{Date}   & \textbf{Milestone} \\
\hline
Mar 29  & Proposal submitted \\
Apr 6   & [e.g., Data collected and preprocessed] \\
Apr 13  & [e.g., Baseline results obtained] \\
Apr 20  & [e.g., Main experiments complete] \\
Apr 27  & [e.g., Draft report finished] \\
[Final] & Report submitted \\
\hline
\end{tabular}
\end{center}


% ─── References ──────────────────────────────────────────────────────────────
\bibliographystyle{plain}
% \bibliography{references}   % uncomment and point to your .bib file

% Manual references example (remove once you switch to BibTeX):
\begin{thebibliography}{9}
  \bibitem{example}
  Author(s).
  \textit{Title of Paper}.
  Venue, Year.
\end{thebibliography}

\end{document}